\section{Desenvolvimento do Backend}

O backend do sistema foi desenvolvido utilizando Node.js como ambiente de execução e Express.js como framework para construção da API REST. A camada de persistência utiliza o Prisma ORM, com banco de dados MySQL, enquanto a autenticação dos usuários é realizada por meio de JSON Web Tokens (JWT), utilizando tokens de acesso e de renovação. Para a validação dos dados recebidos pela API foi utilizado o Zod, e a documentação dos recursos disponibilizados é realizada por meio do padrão Swagger/OpenAPI.

A aplicação foi organizada em módulos de negócio, buscando separar as responsabilidades relacionadas às diferentes funcionalidades do sistema. A estrutura do backend contém módulos específicos para autenticação, usuários, aprendentes, avaliações, aplicações, relatórios, dashboard e auditoria. Além desses módulos, foram definidos diretórios específicos para configurações, middlewares, utilitários e para o mecanismo de auditoria, conforme apresentado na estrutura do projeto.

Essa organização permite concentrar em cada módulo as operações relacionadas a um determinado domínio da aplicação. O módulo de autenticação, por exemplo, concentra as funcionalidades relacionadas ao acesso ao sistema, enquanto os módulos de aprendentes, avaliações e aplicações estão relacionados diretamente ao processo de avaliação psicopedagógica. O módulo de relatórios é responsável pelas operações relacionadas às devolutivas produzidas a partir das aplicações, e o módulo de dashboard disponibiliza informações destinadas à visualização geral dos dados do sistema.

\subsection{Rotas da API}

A comunicação entre o frontend e o backend é realizada por meio de uma API REST, disponibilizada pelo Express.js. Os recursos são organizados em rotas agrupadas por domínio funcional, com controle de acesso aplicado por middleware antes do processamento de cada requisição. A Tabela \ref{qua:rotas_api} apresenta os principais endpoints disponibilizados.

\begin{quadro}[H]
\centering
\small
\caption{Principais rotas da API REST}
\begin{tabular}{p{1.5cm} p{6cm} p{2.5cm} p{4.5cm}}
\hline
\textbf{Método} & \textbf{Rota} & \textbf{Perfil exigido} & \textbf{Descrição} \\
\hline
POST & /api/auth/login & Público & Autenticação e emissão de tokens \\
POST & /api/auth/refresh & Público & Renovação do token de acesso \\
POST & /api/auth/logout & Autenticado & Revogação do refresh token \\
POST & /api/auth/forgot-password & Público & Solicitação de recuperação de senha \\
POST & /api/auth/reset-password & Público & Redefinição de senha via token \\
\hline
GET  & /api/users & ADMIN & Listagem de usuários \\
POST & /api/users & ADMIN & Cadastro de usuário \\
PUT  & /api/users/:id & ADMIN & Atualização de usuário \\
DELETE & /api/users/:id & ADMIN & Exclusão lógica de usuário \\
\hline
GET  & /api/aprendentes & Autenticado & Listagem dos aprendentes do profissional \\
POST & /api/aprendentes & Autenticado & Cadastro de aprendente \\
GET  & /api/aprendentes/:id & Autenticado & Consulta de aprendente \\
PUT  & /api/aprendentes/:id & Autenticado & Atualização de aprendente \\
DELETE & /api/aprendentes/:id & Autenticado & Exclusão lógica de aprendente \\
GET  & /api/aprendentes/:id/historico & Autenticado & Histórico de acompanhamento \\
POST & /api/aprendentes/:id/historico & Autenticado & Inclusão de registro no histórico \\
\hline
GET  & /api/avaliacoes & Autenticado & Listagem dos instrumentos do profissional \\
POST & /api/avaliacoes & Autenticado & Cadastro de instrumento \\
GET  & /api/avaliacoes/:id & Autenticado & Consulta de instrumento \\
PUT  & /api/avaliacoes/:id & Autenticado & Edição (cria nova versão se aplicado) \\
DELETE & /api/avaliacoes/:id & ADMIN & Exclusão de instrumento \\
\hline
GET  & /api/aplicacoes & Autenticado & Listagem de aplicações \\
POST & /api/aplicacoes & Autenticado & Criação de aplicação \\
GET  & /api/aplicacoes/:id & Autenticado & Consulta de aplicação \\
DELETE & /api/aplicacoes/:id & Autenticado & Exclusão (somente se não COMPLETED) \\
PUT  & /api/aplicacoes/:id/answers & Autenticado & Registro ou atualização de respostas \\
POST & /api/aplicacoes/:id/finalize & Autenticado & Finalização da aplicação \\
\hline
GET  & /api/relatorios & Autenticado & Listagem de relatórios \\
GET  & /api/relatorios/:id & Autenticado & Consulta de relatório \\
GET  & /api/relatorios/aplicacao/:id & Autenticado & Relatório por aplicação \\
POST & /api/relatorios & Autenticado & Criação de relatório \\
PUT  & /api/relatorios/:id & Autenticado & Atualização de relatório \\
\hline
GET  & /api/dashboard & Autenticado & Dados do painel geral \\
GET  & /api/auditoria & Autenticado & Consulta do registro de auditoria \\
\hline
\end{tabular}
\fonte{Elaborado pela autora (2026).}
\label{qua:rotas_api}
\end{quadro}

\subsection{Organização dos módulos e middlewares}

A estrutura interna da aplicação está dividida entre configurações, middlewares, módulos de negócio e utilitários. O diretório \textit{config} concentra configurações relacionadas ao ambiente, ao Prisma e ao Swagger. Já o diretório \textit{middlewares} contém mecanismos relacionados à autenticação, validação e tratamento de erros.

Os middlewares são utilizados como uma camada intermediária no processamento das requisições. Dessa forma, a autenticação pode ser aplicada antes do acesso a recursos protegidos, os dados recebidos podem ser submetidos às regras de validação e erros ocorridos durante o processamento podem ser tratados de maneira centralizada.

A validação das informações de entrada é realizada com o Zod. Essa abordagem permite definir as estruturas esperadas para os dados recebidos pela API e verificar se as informações fornecidas estão de acordo com essas estruturas antes de serem processadas pela aplicação.

\subsection{Autenticação e controle de acesso}

A autenticação foi implementada utilizando JWT, com tokens de acesso (\textit{access token}) e de renovação (\textit{refresh token}). O access token tem validade de 15 minutos; o refresh token, de 7 dias. Ambos os valores são configuráveis por variável de ambiente.

No frontend, os tokens são armazenados no \textit{localStorage} do navegador. Essa escolha simplifica a implementação, mas expõe os tokens a ataques de \textit{cross-site scripting} (XSS); a alternativa mais segura seria o uso de cookies \textit{httpOnly}, que não são acessíveis por JavaScript. Essa é uma limitação do protótipo para uso em produção.

O refresh token é rotacionado a cada uso: ao ser utilizado, o token anterior é revogado no banco de dados e um novo é emitido. O token é armazenado em texto plano na tabela \textit{refresh\_tokens}; não é guardado como hash, o que constitui outra limitação de segurança reconhecida. No logout, o token é revogado imediatamente. Na redefinição de senha, todos os refresh tokens ativos do usuário são revogados em uma única transação, invalidando sessões abertas em outros dispositivos.

O processo de recuperação de senha utiliza um token gerado com \texttt{crypto.randomBytes(32)}, com validade de 1 hora. O link de redefinição é enviado por e-mail por meio do Nodemailer, configurado com servidor SMTP definido em variáveis de ambiente. O token de redefinição é marcado como utilizado após o uso, impedindo sua reutilização.

Para proteção contra ataques de força bruta e abuso da API, o backend aplica limitação de taxa (\textit{rate limiting}) globalmente por meio da biblioteca \textit{express-rate-limit}: no máximo 100 requisições por endereço IP em uma janela de 15 minutos. Ao atingir o limite, a API retorna HTTP 429 com mensagem de erro. Os valores são configuráveis por variável de ambiente.

O backend contempla os perfis \textit{ADMIN} e \textit{PSICOPEDAGOGO}, com controle de acesso por perfil aplicado via middleware \textit{authorize} nas rotas que exigem permissão específica. Além disso, o sistema implementa restrição de visibilidade baseada em titularidade: aprendentes e avaliações são vinculados ao usuário que os criou, e aplicações e relatórios seguem a titularidade do aprendente a que pertencem. As operações de leitura e escrita alcançam somente os registros pertencentes ao usuário autenticado, exceto para o perfil \textit{ADMIN}, que possui acesso irrestrito.

\subsection{Gerenciamento de aprendentes e avaliações}

O módulo de aprendentes disponibiliza as operações necessárias para o cadastro e gerenciamento dos indivíduos acompanhados pelo profissional. Os dados são posteriormente utilizados nas aplicações dos instrumentos avaliativos, mantendo a associação entre o aprendente e os registros produzidos durante o processo de avaliação.

A visibilidade dos aprendentes é controlada pelo campo \textit{responsavelId}, preenchido automaticamente com o identificador do usuário autenticado no momento do cadastro. Nas operações de listagem e consulta individual, o módulo aplica um filtro que restringe os resultados ao usuário responsável. O perfil \textit{ADMIN} constitui exceção a essa regra, tendo acesso irrestrito a todos os registros. Essa lógica é implementada na camada de serviço, antes da consulta ao banco de dados, sem expor ao cliente a distinção entre registros próprios e de outros usuários.

O módulo de avaliações é responsável pelo gerenciamento dos instrumentos utilizados pelo sistema. A estrutura adotada permite representar avaliações compostas por diferentes questões e tipos de resposta. O versionamento desses instrumentos é mantido na camada de persistência, conforme apresentado na modelagem do banco de dados, permitindo relacionar diferentes versões de uma mesma avaliação.

De maneira análoga ao módulo de aprendentes, a visibilidade das avaliações é controlada pelo campo \textit{autorId}, atribuído automaticamente ao usuário que criou o instrumento. As consultas de listagem e de detalhe aplicam o mesmo filtro de restrição por responsável, garantindo que cada profissional acesse somente os instrumentos que criou. Ao criar uma nova versão de uma avaliação, o campo \textit{autorId} recebe o usuário que realizou a alteração. Como o psicopedagogo só edita os próprios instrumentos, a titularidade se mantém entre as versões; quando a edição é feita pelo administrador, a nova versão passa a pertencer a ele, limitação discutida no Capítulo \ref{capitulo10}.

As aplicações são tratadas separadamente pelo módulo \textit{aplicacoes}, permitindo associar uma avaliação a um aprendente e registrar o andamento desse processo. Essa separação entre avaliação e aplicação possibilita que um mesmo instrumento seja utilizado em diferentes processos avaliativos, mantendo cada aplicação como um registro independente.

\subsection{Relatórios e dashboard}

O módulo \textit{relatorios} concentra as operações relacionadas aos relatórios produzidos a partir das aplicações realizadas. Esses registros permanecem associados à aplicação correspondente, permitindo relacionar a devolutiva produzida pelo profissional ao processo avaliativo que lhe deu origem.

O backend também possui um módulo específico de \textit{dashboard}, disponibilizado por meio da rota \textit{GET /api/dashboard}. Esse módulo fornece os dados necessários para a apresentação das informações gerais do sistema na interface, evitando que o frontend precise realizar diretamente operações sobre o banco de dados.

Nesse contexto, o backend atua como intermediário entre os dados persistidos e a interface do sistema, sendo responsável pelo processamento das requisições e pela disponibilização das informações necessárias às funcionalidades implementadas.

\subsection{Documentação e testes}

A API possui documentação baseada em Swagger/OpenAPI, disponibilizada durante a execução do sistema no endereço \textit{/api-docs}. A documentação permite consultar os recursos disponibilizados pela API e auxilia na compreensão e utilização dos endpoints durante o desenvolvimento e os testes.

Para os testes automatizados do backend foram utilizadas as ferramentas Vitest e Supertest. O Vitest é responsável pela execução e pela geração do relatório de cobertura, enquanto o Supertest possibilita a realização de requisições HTTP reais contra a aplicação em memória, sem necessidade de servidor externo. O frontend utiliza Vitest combinado com Testing Library para testes de componentes React.

O backend conta com 76 testes automatizados distribuídos em dois níveis: testes de integração, que exercitam as rotas da API de ponta a ponta, e testes unitários, voltados para as regras de negócio isoladas dos módulos de autenticação e do mecanismo de auditoria. A Tabela~\ref{qua:testes_backend} apresenta a distribuição por arquivo de teste e a cobertura de linhas por módulo obtida com o relatório gerado pelo Vitest com provider V8.

\begin{quadro}[H]
\centering
\small
\caption{Distribuição dos testes e cobertura por módulo do backend}
\begin{tabular}{p{5cm} p{2.2cm} p{2.5cm} p{2.5cm}}
\hline
\textbf{Módulo / Arquivo de teste} & \textbf{Nível} & \textbf{Testes} & \textbf{Cobertura (linhas)} \\
\hline
auth (rotas) - auth.test.js & Integração & 7 & 68,32\% \\
aprendentes - aprendentes.test.js & Integração & 6 & 74,85\% \\
avaliacoes - avaliacoes.test.js & Integração & 10 & 78,57\% \\
aplicacoes - aplicacoes.test.js & Integração & 14 & 95,83\% \\
relatorios - relatorios.test.js & Integração & 8 & 95,67\% \\
auditoria - auditoria.test.js & Integração & 8 & 97,08\% \\
auth.service - auth.service.test.js & Unitário & 4 & - \\
Extensão Prisma - audit.extension.test.js & Unitário & 12 & 95,37\% \\
Utilitários - shared-utils.test.js & Unitário & 7 & 70,14\% \\
\hline
\textbf{Total} & & \textbf{76} & \textbf{80,20\% (linhas)} \\
 & & & \textbf{70,27\% (branches)} \\
\hline
\end{tabular}
\fonte{Elaborado pela autora (2026).}
\label{qua:testes_backend}
\end{quadro}

A cobertura de linhas de 80,20\% e de branches de 70,27\% superam o critério mínimo de 70\% estabelecido em RNF10. Os módulos com menor cobertura são \textit{users} (64,31\%) e \textit{auth.service} (42,50\% de linhas nos testes unitários isolados), reflexo de operações administrativas e de recuperação de senha que não possuem testes de integração completos.

O frontend conta com 51 testes distribuídos em 9 arquivos, cobrindo as páginas e componentes de maior complexidade de lógica: LoginPage, ProtectedRoute, AvaliacaoFormPage, AvaliacaoDetailPage, AplicacaoDetailPage, Auditoria, RelatorioDetailPage, StatusChip e a função utilitária de impressão. A cobertura do frontend atingiu 50,44\%.

Também foram definidos procedimentos para geração do cliente do Prisma, execução das migrações e população inicial do banco de dados. Essas operações permitem preparar o ambiente de desenvolvimento de forma padronizada e manter a estrutura do banco sincronizada com o modelo utilizado pela aplicação.

\subsection{Auditoria}

A rastreabilidade das operações é garantida por um registro de auditoria, armazenado na entidade \textit{audit\_logs}. As razões que motivaram esse registro, relacionadas às lacunas dos campos de controle do modelo de dados, são discutidas no Capítulo \ref{capitulo13}; esta seção descreve sua implementação.

O preenchimento é feito por uma extensão do Prisma Client, mecanismo que substitui os antigos \textit{middlewares} do Prisma: todas as operações executadas pela aplicação sobre o banco de dados passam por essa extensão antes de serem concluídas. Para que a extensão saiba qual usuário originou cada operação, o backend mantém um contexto por requisição, implementado com o \textit{AsyncLocalStorage} do Node.js. O contexto é criado no início do processamento de cada requisição, com o endereço IP de origem, e recebe o identificador e o perfil do usuário quando o middleware de autenticação valida o token de acesso. Dessa forma, os serviços de negócio não precisam repassar essas informações à camada de persistência.

Nas operações de criação, alteração e exclusão, a extensão grava a ação realizada, a entidade, o identificador do registro, o usuário, o perfil, o IP e o momento da operação. Nas alterações e exclusões, os valores anteriores são lidos antes da operação, e o registro guarda somente os campos alterados, com seus valores anteriores e novos. As exclusões lógicas, feitas pelo preenchimento do campo \textit{deletedAt}, são registradas como exclusões. O hash da senha nunca é gravado: quando ele muda, o registro indica apenas que o campo foi alterado.

Nas entidades com dados clínicos de crianças (aprendentes, histórico de acompanhamento, aplicações e relatórios), também as consultas são registradas, inclusive as listagens, com um registro para cada item exibido. Cada registro é associado ao aprendente a que o dado pertence, o que permite responder quem acessou os dados de uma determinada criança. Os logins bem-sucedidos e as tentativas recusadas também são registrados, estas sem a senha informada.

A consulta ao registro é feita pelo endpoint \textit{GET /api/auditoria}, com filtros por aprendente, usuário, entidade, ação e período. O perfil \textit{ADMIN} consulta todo o registro. O perfil \textit{PSICOPEDAGOGO} consulta apenas o registro de acessos dos aprendentes sob sua responsabilidade, o que inclui os acessos feitos pelo administrador aos dados dessas crianças.

As limitações desse registro de auditoria são discutidas no Capítulo \ref{capitulo13}.

Dessa forma, o desenvolvimento do backend estabeleceu a camada responsável pela disponibilização da API REST, autenticação, validação das requisições, aplicação das operações relacionadas aos módulos do sistema, registro de auditoria, comunicação com o banco de dados e disponibilização dos dados utilizados pelo frontend. A organização modular adotada também fornece uma estrutura para a manutenção e evolução das funcionalidades relacionadas ao processo de avaliação psicopedagógica.